% Source: https://artofproblemsolving.com/community/c3h1199712_emcc_problem
% Stable identity: aops-topic-1199712-q01
Let the side be $s$ and put $P=(0,0)$, $Q=(s,0)$, $S=(0,s)$. Then $M=\left(\dfrac {s}{8},0\right)$ and $N=\left(0,\dfrac {s}{8}\right)$. Line $MS$ is $8x+y=s$ and line $NQ$ is $x+8y=s$, meeting at $X=\left(\dfrac {s}{9},\dfrac {s}{9}\right)$.

\begin{asy}
usepackage("amsmath");
size(3.5cm);
defaultpen(fontsize(10pt));
real s=1; pair P=(0,0),Q=(s,0),R=(s,s),S0=(0,s),M=(s/8,0),N=(0,s/8),X=(s/9,s/9);
draw(P--Q--R--S0--cycle,linewidth(0.9)); draw(M--S0,blue); draw(N--Q,blue); fill(P--M--X--N--cycle,red);
label("$P$",P,SW);label("$Q$",Q,SE);label("$R$",R,NE);label("$S$",S0,NW);label("$M$",M,S);label("$N$",N,W);label("$X$",X,NE);
label("\mbox{}",(min(currentpicture,user=true)+max(currentpicture,user=true))/2,black); // keeps the text after the figure black when LaTeX typesets the labels inline
\end{asy}

By symmetry in the line $y=x$, $PMXN$ consists of two congruent triangles $PMX$ and $PNX$, each of area $\dfrac{1}{2}\cdot\dfrac {s}{8}\cdot\dfrac {s}{9}=\dfrac{s^2}{144}$. So the area is $\dfrac{s^2}{72}=\dfrac{1}{72}$ and $s=\boxed{1}$.
